\chapter{Các thuật toán khác và Thủ thuật}

\section{Xử lý đoạn (Intervals)}
	\kactlimport{IntervalContainer.h}
	\kactlimport{IntervalCover.h}
	\kactlimport{ConstantIntervals.h}

\section{Các thuật toán khác}
	\kactlimport{TernarySearch.h}
	\kactlimport{LIS.h}
	\kactlimport{FastKnapsack.h}

\section{Quy hoạch động nâng cao}
	\kactlimport{KnuthDP.h}
	\kactlimport{DivideAndConquerDP.h}

\section{Mẹo gỡ lỗi (Debugging)}
	\begin{itemize}
		\item \verb@signal(SIGSEGV, [](int) { _Exit(0); });@ chuyển Runtime Error (segfault) thành Wrong Answer.
			Tương tự có thể bắt SIGABRT (assertion thất bại) và SIGFPE (chia cho 0).
		\item \verb@feenableexcept(29);@ dừng chương trình ngay khi gặp NaN (\texttt 1), chia cho 0 (\texttt 4), vô cùng (\texttt 8) và denormals (\texttt{16}).
	\end{itemize}

\section{Tối ưu tốc độ (Optimization)}
	\verb@__builtin_ia32_ldmxcsr(40896);@ tắt denormals (nguyên nhân khiến số thực chậm gấp 20 lần khi gần về 0).
	\subsection{Thao tác bit (Bit hacks)}
		\begin{itemize}
			\item \verb@x & -x@ là bit 1 nhỏ nhất (bật thấp nhất) trong \texttt{x}.
			\item \verb@for (int x = m; x; ) { --x &= m; ... }@ duyệt qua tất cả mặt nạ tập con của \texttt{m} (ngoại trừ chính \texttt{m}).
			\item \verb@c = x&-x, r = x+c; (((r^x) >> 2)/c) | r@ tìm số tiếp theo lớn hơn \texttt{x} có cùng số lượng bit 1.
			\item \verb@rep(b,0,K) rep(i,0,(1 << K))@ \\ \verb@  if (i & 1 << b) D[i] += D[i^(1 << b)];@ tính tổng trên mọi tập con (SOS DP).
		\end{itemize}
	\subsection{Chỉ thị Pragma}
		\begin{itemize}
			\item \lstinline{#pragma GCC optimize ("Ofast")} tự động vector hóa vòng lặp và tối ưu số thực.
			\item \lstinline{#pragma GCC target ("avx2")} tăng gấp đôi hiệu năng nhờ vectorization hiện đại.
			\item \lstinline{#pragma GCC optimize ("trapv")} dừng chương trình khi tràn số nguyên có dấu (chạy khá chậm).
		\end{itemize}
	\kactlimport{FastMod.h}
	\kactlimport{FastInput.h}
	\kactlimport{BumpAllocator.h}
	\kactlimport{SmallPtr.h}
	\kactlimport{BumpAllocatorSTL.h}
	% \kactlimport{Unrolling.h}
	\kactlimport{SIMD.h}
